% ============================================================
%  00 — Front Matter : Page de garde, Dédicaces, Remerciements,
%       Résumés (FR / EN / AR), TOC, Liste des figures, Abréviations
% ============================================================

% ─── Page de garde ───────────────────────────────────────────
\begin{titlepage}
\thispagestyle{empty}
\begin{center}

% En-tête institutionnel
{\footnotesize\textbf{République Tunisienne}}\\[2pt]
{\footnotesize Ministère de l'Enseignement Supérieur et de la Recherche Scientifique}\\[12pt]

\begin{minipage}[c]{0.45\textwidth}
  \centering
  \begin{tcolorbox}[enhanced, colback=farmblue!10, colframe=farmblue,
    arc=6pt, boxrule=1pt, width=3.5cm, height=2cm, valign=center, halign=center]
    {\small\bfseries\color{farmblue} Tek-Up\\École d'Ingénieurs}
  \end{tcolorbox}
\end{minipage}%
\begin{minipage}[c]{0.45\textwidth}
  \centering
  \begin{tcolorbox}[enhanced, colback=farmgreen!10, colframe=farmgreen,
    arc=6pt, boxrule=1pt, width=3.5cm, height=2cm, valign=center, halign=center]
    {\small\bfseries\color{farmgreen} PFE 2024-2025\\Data Science \& IA}
  \end{tcolorbox}
\end{minipage}

\vspace{0.4cm}
\rule{\textwidth}{1.5pt}\\[4pt]
{\large\textbf{École Privée d'Ingénieurs — Tek-Up}}\\
{\normalsize Spécialité : \textit{Data Science et Intelligence Artificielle}}\\[4pt]
\rule{\textwidth}{1.5pt}

\vspace{0.8cm}
{\large\textsc{Rapport de Fin de Cycle d'Ingénierie}}\\[6pt]
{\normalsize En vue de l'obtention du diplôme d'\\[2pt]
\textbf{Ingénieur en Data Science et Intelligence Artificielle}}

\vspace{0.8cm}

% Titre du projet
\begin{tcolorbox}[
  enhanced, colback=farmblue!8, colframe=farmblue,
  arc=8pt, boxrule=1.5pt,
  width=\textwidth
]
\centering
{\LARGE\bfseries\color{farmblue} Smart Farm AI v3.0}\\[6pt]
{\large\color{farmgreen}
  Plateforme Agro-Technologique Multi-Tenant basée sur\\
  l'Intelligence Artificielle, l'IoT et les LLM Souverains}\\[4pt]
{\normalsize\color{farmgray}\textit{Détection automatisée des maladies agricoles, gestion apicole intelligente\\
et assistant conversationnel en Darija tunisien}}
\end{tcolorbox}

\vspace{0.8cm}

% Auteur
\begin{minipage}{0.5\textwidth}
  \raggedright
  {\normalsize\textbf{Réalisé par :}}\\[3pt]
  {\large\color{farmblue}\textbf{Mohamed Sayari}}\\
  {\small Classe : DSI-5 (Data Science 5\textsuperscript{ème} année)}
\end{minipage}%
\begin{minipage}{0.5\textwidth}
  \raggedleft
  {\normalsize\textbf{Encadré par :}}\\[3pt]
  {\large\textbf{[Nom de l'encadrant]}}\\
  {\small Maître de conférences, Tek-Up}
\end{minipage}

\vspace{0.6cm}

% Jury
\begin{tcolorbox}[
  enhanced, colback=gray!5, colframe=farmgray,
  arc=4pt, boxrule=0.8pt, width=0.88\textwidth
]
\centering
\textbf{Jury de soutenance}\\[4pt]
\begin{tabular}{ll}
Président du jury  & [Nom Prénom], Professeur, [Institution]\\
Rapporteur         & [Nom Prénom], Maître de conférences, [Institution]\\
Encadrant          & [Nom Prénom], Maître de conférences, Tek-Up\\
\end{tabular}
\end{tcolorbox}

\vfill
{\normalsize Année universitaire : \textbf{2024--2025}}

\end{center}
\end{titlepage}

% ─── Page vierge après couverture ────────────────────────────
\cleardoublepage

% ─── Dédicaces ───────────────────────────────────────────────
\chapter*{Dédicaces}
\thispagestyle{empty}
\addcontentsline{toc}{chapter}{Dédicaces}

\begin{center}
\vspace{2cm}
{\large\itshape
À mes parents, pour leur soutien indéfectible\\
et leurs sacrifices qui ont rendu ce parcours possible.\\[1.5em]
À mes professeurs de Tek-Up,\\
qui m'ont transmis la passion pour la Data Science et l'IA.\\[1.5em]
À tous les agriculteurs tunisiens\\
dont le travail quotidien a inspiré ce projet.
}
\vspace{2cm}

\begin{flushright}
\textit{Mohamed Sayari}
\end{flushright}
\end{center}

\cleardoublepage

% ─── Remerciements ────────────────────────────────────────────
\chapter*{Remerciements}
\thispagestyle{empty}
\addcontentsline{toc}{chapter}{Remerciements}

Je tiens à exprimer ma profonde gratitude à toutes les personnes qui ont contribué, de près ou de loin, à la réalisation de ce projet de fin de cycle.

Mes sincères remerciements vont en premier lieu à mon encadrant pédagogique, dont les conseils avisés, la disponibilité et la rigueur scientifique ont guidé chaque étape de ce travail.

Je remercie également le corps enseignant de l'École Tek-Up pour la formation de qualité dispensée tout au long de ces cinq années d'ingénierie en Data Science et Intelligence Artificielle.

Ma gratitude s'adresse aussi à l'Union Tunisienne de l'Agriculture et de la Pêche (UTAP) et à l'Agence de Vulgarisation et de Formation Agricoles (AVFA) pour les ressources documentaires sur les pratiques agricoles et apicoles tunisiennes, qui ont enrichi le corpus RAG de l'assistant souverain.

Enfin, je remercie ma famille pour son soutien moral constant et sa confiance tout au long de ce parcours académique.

\cleardoublepage

% ─── Résumé Français ─────────────────────────────────────────
\chapter*{Résumé}
\thispagestyle{empty}
\addcontentsline{toc}{chapter}{Résumé}

\noindent\textbf{Titre :} Smart Farm AI v3.0 — Plateforme Agro-Technologique Multi-Tenant basée sur l'IA et l'IoT

\medskip
Ce rapport présente le développement de \farmAI{} v3.0, une plateforme numérique innovante dédiée à l'agriculture de précision en Tunisie. Face aux défis de la productivité agricole — contraintes hydriques, ravageurs, maladies des cultures et déclin des populations de pollinisateurs — ce projet propose une solution intégrée combinant l'Internet des Objets (IoT), la vision par ordinateur et les grands modèles de langage (LLM).

L'architecture repose sur un backend FastAPI avec plus de 50 endpoints REST, un frontend React/Vite responsive, et deux nœuds ESP32 pour la surveillance en temps réel : \textbf{NODE\_A} pour l'irrigation intelligente (détection d'humidité, pression de canalisation, débit) et \textbf{NODE\_B} pour la gestion apicole (température de ruche, poids, humidité extérieure).

Le module de vision artificielle intègre huit modèles YOLO v11 couvrant treize espèces et pathologies : abeilles (détection de ruche par boîtes orientées), maladies foliaires (12 classes), maladies du bétail, du troupeau ovin/caprin, de l'olivie, des agrumes, des insectes ravageurs et la détection incendie/fumée. La précision moyenne (mAP50) atteint 91,4\% sur le modèle de détection des ruches.

L'assistant conversationnel souverain utilise le modèle Labess-7B — un LLM spécialisé dans le dialecte arabe tunisien (Darija) — enrichi d'un système RAG (Retrieval-Augmented Generation) alimenté par ChromaDB et des sources officielles tunisiennes (UTAP, AVFA). Il intègre la reconnaissance vocale, la synthèse vocale et l'analyse d'images multimodale via LLaVA.

La chaîne MLOps s'appuie sur DVC pour le versionnement des données et pipelines, et MLflow pour le suivi des expériences. Un modèle IsolationForest détecte les anomalies télémétriques (score de silhouette : 0,73) et déclenche des alertes en temps réel.

Le projet a été conduit selon la méthodologie Scrum en 6 sprints sur 14 semaines, complétée par le cycle CRISP-DM pour la partie Data Science. La plateforme est déployée sur Oracle Cloud Free Tier (ARM64) avec Docker Compose.

\medskip
\noindent\textbf{Mots-clés :} Agriculture de précision, YOLO, IoT, ESP32, RAG, Darija, MLOps, FastAPI, React, Apiculture, Tunisie.

\cleardoublepage

% ─── Résumé Anglais ──────────────────────────────────────────
\chapter*{Abstract}
\thispagestyle{empty}
\addcontentsline{toc}{chapter}{Abstract}

\noindent\textbf{Title:} Smart Farm AI v3.0 — Multi-Tenant Agro-Technological Platform based on AI and IoT

\medskip
This report presents the development of \farmAI{} v3.0, an innovative digital platform for precision agriculture in Tunisia. Facing agricultural productivity challenges — water constraints, pests, crop diseases, and declining pollinator populations — this project proposes an integrated solution combining Internet of Things (IoT), computer vision, and Large Language Models (LLMs).

The architecture is built on a FastAPI backend with 50+ REST endpoints, a responsive React/Vite frontend, and two ESP32 nodes for real-time monitoring: \textbf{NODE\_A} for smart irrigation (humidity detection, pipe pressure, flow rate) and \textbf{NODE\_B} for beehive management (hive temperature, weight, external humidity).

The computer vision module integrates eight YOLO v11 models covering thirteen species and pathologies: bees (oriented bounding box detection), leaf diseases (12 classes), livestock diseases, sheep/goat, olive, citrus, pest insects, and fire/smoke detection. The average precision (mAP50) reaches 91.4\% on the beehive detection model.

The sovereign conversational assistant uses the Labess-7B model — an LLM specialized in the Tunisian Arabic dialect (Darija) — enhanced with a RAG (Retrieval-Augmented Generation) system powered by ChromaDB and official Tunisian sources (UTAP, AVFA). It integrates speech recognition, text-to-speech, and multimodal image analysis via LLaVA.

The MLOps pipeline relies on DVC for data and pipeline versioning, and MLflow for experiment tracking. An IsolationForest model detects telemetric anomalies (silhouette score: 0.73) and triggers real-time alerts.

The project was conducted using Scrum methodology over 6 sprints across 14 weeks, complemented by the CRISP-DM cycle for the Data Science component. The platform is deployed on Oracle Cloud Free Tier (ARM64) with Docker Compose.

\medskip
\noindent\textbf{Keywords:} Precision agriculture, YOLO, IoT, ESP32, RAG, Darija, MLOps, FastAPI, React, Beekeeping, Tunisia.

\cleardoublepage

% ─── Table des matières ───────────────────────────────────────
\tableofcontents
\cleardoublepage

% ─── Liste des figures ────────────────────────────────────────
\listoffigures
\addcontentsline{toc}{chapter}{Liste des Figures}
\cleardoublepage

% ─── Liste des tableaux ───────────────────────────────────────
\listoftables
\addcontentsline{toc}{chapter}{Liste des Tableaux}
\cleardoublepage

% ─── Liste des abréviations ───────────────────────────────────
\chapter*{Liste des Abréviations}
\addcontentsline{toc}{chapter}{Liste des Abréviations}
\thispagestyle{empty}

\begin{longtable}{@{}lp{10cm}@{}}
\toprule
\textbf{Abréviation} & \textbf{Signification} \\
\midrule
\endhead
\bottomrule
\endfoot
API     & Application Programming Interface \\
AVFA    & Agence de Vulgarisation et de Formation Agricoles \\
CORS    & Cross-Origin Resource Sharing \\
CPU     & Central Processing Unit \\
CRISP-DM & Cross-Industry Standard Process for Data Mining \\
CSV     & Comma-Separated Values \\
CV      & Computer Vision (Vision par Ordinateur) \\
DVC     & Data Version Control \\
ESP32   & Microcontrôleur Wi-Fi/Bluetooth d'Espressif Systems \\
FAO     & Food and Agriculture Organization (ONU) \\
GPU     & Graphics Processing Unit \\
HTTP    & HyperText Transfer Protocol \\
IA      & Intelligence Artificielle \\
INS     & Institut National de la Statistique \\
IoT     & Internet of Things (Internet des Objets) \\
JSON    & JavaScript Object Notation \\
JWT     & JSON Web Token \\
LLM     & Large Language Model (Grand Modèle de Langage) \\
mAP     & Mean Average Precision \\
MLOps   & Machine Learning Operations \\
MQTT    & Message Queuing Telemetry Transport \\
MVP     & Minimum Viable Product \\
NLP     & Natural Language Processing \\
OBB     & Oriented Bounding Box (Boîte Englobante Orientée) \\
OCR     & Optical Character Recognition \\
ORM     & Object-Relational Mapping \\
OTP     & One-Time Password \\
PIB     & Produit Intérieur Brut \\
PWA     & Progressive Web App \\
RAG     & Retrieval-Augmented Generation \\
REST    & Representational State Transfer \\
RGB     & Red-Green-Blue \\
SQL     & Structured Query Language \\
STT     & Speech-to-Text \\
SIG/GIS & Système d'Information Géographique / Geographic Information System \\
TTS     & Text-to-Speech \\
UTAP    & Union Tunisienne de l'Agriculture et de la Pêche \\
VPN     & Virtual Private Network \\
WSGI    & Web Server Gateway Interface \\
YOLO    & You Only Look Once \\
\bottomrule
\end{longtable}

\cleardoublepage
